\section{DLA Case-Study Deployment}
\label{sec:deployment}

\subsection{Static Engine Interface}

The model is fused and exported to ONNX with a fixed batch size and input
resolution. TensorRT parses the graph in FP16 and targets DLA core~0 of the
module; the engine exposes one image input and the three stride-ordered
prediction bindings of Sec.~\ref{sec:retrofit}. Because the accelerator is the
default device and fallback is disabled, the compiler either accepts the whole
graph or refuses to build, which is the property we use as an acceptance test.

\subsection{Confidence-First Sparse DFL Decode}

The host decoder splits the first $4r=64$ channels from the $n_c=80$ class
logits at each scale. Since a candidate's class confidence does not depend on
its box distribution, the decoder first discards every anchor whose maximum
class logit falls below the inverse-sigmoid image of the confidence threshold,
$\log\frac{\tau}{1-\tau}$, and only then restores grids and applies the 16-bin
expectation \eqref{eq:dfl} to the survivors, converts distances to boxes,
applies $\sigma(\cdot)$, and runs standard NMS. Up to floating-point roundoff
this is equivalent to dense decoding followed by the same threshold, because
thresholding commutes with a per-anchor box transform that it does not read;
the dense path is retained in the implementation as an equivalence gate and is
the reference against which the sparse path is checked. The saving is large
because the surviving fraction is small: at $\tau=0.25$ on the test image, 13 of
$8400$ anchors survive (Sec.~\ref{sec:rq3}).

\subsection{DLA-Native Input Packing and Output Unpacking}
\label{sec:packing}

The strict engine does not accept a linear NCHW image and does not return
linear NCHW predictions, and this is the single most consequential difference
between deploying on the accelerator and deploying on the GPU. It is built with
\texttt{--inputIOFormats=fp16:dla\_hwc4} and
\texttt{--outputIOFormats=fp16:chw16}, and the build-time layer dump confirms
the input as the DLA image format, ``channel major FP16 where
$C \bmod 4 = 0$'', and the three outputs as ``sixteen wide channel vectorized
row major FP16''. At $640\times640$ the input is a 3-channel image padded to 4
channels and interleaved per pixel, with each row aligned to 64 bytes, giving a
$4\times640\times640$ FP16 buffer of $3.28$~MB; the outputs need no channel
padding, since $144=9\times16$, and total $144\times8400$ FP16 values, i.e.\
$2.42$~MB. \TODO{the runtime-reported binding format string for this engine is
not recorded in our logs; the TensorRT Python API may report a
channel-vectorized alias for the DLA image format, which changes nothing about
the buffer geometry but should be recorded in the final runs}

Where this conversion runs, and in what implementation, is a property of the
integrator and not of the accelerator, so we state it exactly. In our reference
pipeline both directions run \emph{on the host CPU in single-threaded PyTorch}
(\texttt{torch.set\_num\_threads(1)}). Packing allocates a zero-filled padded
tensor, copies a permuted NCHW-to-interleaved view into it, flattens it, and
copies it synchronously into the device buffer; unpacking copies the device
buffer to the host and then applies a view, a permutation, a
\texttt{contiguous()} materialization, and a float cast per output. There is no
CUDA kernel, no pinned staging buffer, no NEON-tuned or vectorized packer, and
no overlap between conversion and execution. This is a deliberately simple
reference implementation, and Sec.~\ref{sec:rq3} reports its cost as such.

\subsection{Compliance Gates}

We instantiate $A_{\mathcal{H}}(G)=1$ as strict DLA execution only if all of the
following hold, in order:
\begin{enumerate}
  \item the model is fused before export, and every engine input and learned
  output has a fixed four-dimensional shape;
  \item the exported ONNX graph contains no rank-changing terminal subgraph, no
  matrix multiplication, and no dynamic shape operator, audited on the graph
  rather than assumed from the exporter flags;
  \item TensorRT builds with DLA as the default device and GPU fallback
  disabled, so that a single inadmissible node is a build failure;
  \item the verbose build log and the engine inspector report no GPU-assigned
  learned layer, apart from unavoidable host-side I/O handling; and
  \item externally decoded predictions agree with the ordinary PyTorch
  inference path within a documented tolerance, and the sparse decoder agrees
  with the dense one.
\end{enumerate}
Gate~2 is an audit of the graph the compiler will see, not a whitelist claim
about the source model: the retrofit inherits a static channel split, a channel
softmax, and SiLU decompositions whose fusion is a compiler behavior that can
change across toolchain versions, so gate~3 remains the binding criterion and
gate~2 exists to explain build failures rather than to predict success. Gates 3
and 4 are evaluated on the target module, because engines and DLA loadables are
specific to the hardware and software stack that produced them. A
fallback-enabled build never satisfies these gates and is reported only as a
diagnostic. On a different accelerator, gates 1, 3, 4, and 5 carry over
verbatim and only the operator audit of gate~2 is replaced.
